\clearpage
\fancyhead[L]{\small ROHDE. ON THE MASSES OF THE PLANETS.}
\SourceLeaf{488}
\begin{center}
{\large ROHDE. ON THE MASSES OF THE PLANETS.\par}
\medskip
\rule{0.55\linewidth}{0.4pt}\\[2pt]
von \textsc{Zach}. Monatliche Correspondenz. Vol. XII. pp. 413 . . 419. October 1805.\\[2pt]
\rule{0.55\linewidth}{0.4pt}
\end{center}
\begin{center}Opinion of another reviewer on Captain \textsc{Rohde}’s paper on the masses of the planets, taken from a letter by that reviewer.\end{center}

\noindent --- --- --- The --- --- --- defence by Captain \textsc{Rohde} that you kindly sent me amused me somewhat as an example of an even worse defence of a bad cause. Still, I do regret that I am partly to blame for starting such a thankless dispute. My hope that a rebuke would at least make Captain \textsc{Rohde} more cautious has come to nothing. His ill-considered reply, which he had printed in such haste, is once again full of fresh proofs of his --- --- --- --- --- I take up the matter with you once more with no wish to anticipate your answer (you have probably already finished it, and the thoroughness and deep conviction that shine through your review are a guarantee that no one can defend it better than you yourself). My main reason is that you want to know what I think of this reply to your criticism.

Here is my view of the dispute. In your review you left open the possibility that \textsc{Rohde}’s method of determining the masses of the planets might be correct and reliable in the abstract, or on purely theoretical grounds. You asserted and proved

1) that \textsc{Rohde} applied it quite wrongly to Ceres and Pallas, so that for that reason alone his results for them cannot be worth anything at all.

2) that with practical astronomy in its present state, it is impossible to determine the mean distances of the planets as accurately as this method would require unless they are calculated from the orbital period with their masses taken into account.

I have never come across a defence like Captain \textsc{Rohde}’s against 1). In other words, it amounts to this. \emph{‘Since my method, as the reviewer has proved in 2), is quite useless and of no practical value after all, it would have been ridiculous to take any special care in applying it.’} Unfortunately, in doing so he has made yet another schoolboy error. He reasons that the precession of the equinoxes in one year amounts to 0.01419 in time, and thus to less than 0.07 days in 1684 days. One would not forgive a beginner in astronomy for such a thing. The precession of the equinoxes in a year is 50$''$, and the Sun takes 0.014 days to catch up with it. The Sun therefore takes 0.07 days to make up\SourceLeaf{489} for the precession of 1684 days, but Ceres takes more than 0.3 days, and that is how much longer the sidereal revolution is than the tropical one.

Everything \textsc{Rohde} brings up against 2) is utterly empty chatter that has nothing to do with the matter. He thinks that the art of observation has advanced so far nowadays that the sixth decimal place of the mean distances can be regarded as reliable. Yet it is well known that all the values given in astronomical works for the mean distances of the planets and of \textsc{Halley}’s comet have been calculated from the orbital periods, either with or without taking the masses into account. So \textsc{Rohde} had no right to use the figures in \textsc{La Place}’s \emph{Exposition du Système du Monde} etc. It is therefore sheer delusion for him to claim that the agreement of his results for planets that have satellites with the results obtained from those figures confirms the validity of his method.

Besides, \textsc{Rohde} seems to have realized himself how worthless his answer was (the contrary would be incomprehensible unless one were prepared to make the strangest assumptions). For he tries to back down and give the impression that he never intended to provide better determinations of the planetary masses. He claims he merely sought a strictly correct method that would withstand every objection, without concerning himself with whether observations could already give the necessary elements and data with sufficient accuracy. There are three remarks to make on this.

\noindent I) \emph{This claim is untrue}, for Mr \textsc{Rohde} asserts quite clearly at the end of his paper that all perturbation calculations so far performed for Mercury, Mars and Venus are worthless, partly because entirely wrong masses were assumed for them and partly because the incomparably large effects of Ceres, Pallas and Juno had to be neglected. He also says that he wanted to tell astronomers about his important discovery by the speedy route of print instead of the slow circulation of a manuscript, so that everyone who could calculate would immediately join forces to recast the perturbation calculations completely using his new values for the masses, \emph{pour accélérer de tous côtés la récolte, à laquelle les forces d'un seul homme ne suffisent pas, et les miennes point du tout} (surely the truest thing in the whole paper), \emph{vû la grande quantité de combinaisons analytiques nouvelles etc.}

\noindent II) Whether one wants to find the mean distance of a planet without calculating it from the orbital period, or to find any distance from observations at all, the only way is to calculate it trigonometrically in the familiar triangle whose three angles are the commutation, elongation and annual parallax, taking the distance of the Sun from the Earth as the known side. The angles are found from the longitude of the Sun and the geocentric and heliocentric longitudes of the planet. The latter must be calculated from the tables. As is well known, however, one cannot vouch for the accuracy of this calculated longitude to within 5$''$ and more for any planet. Nor has the observed geocentric longitude of the Sun and planets been obtained to within a couple of seconds. These figures are only very low estimates. From them one can readily obtain results for the various planets that show that the resulting errors in the distances are enormous compared with the accuracy required, and make \textsc{Rohde}’s method quite useless.

Once \textsc{Rohde} is driven out of this hiding place, he will content himself with saying that his method is correct in theory after all, and that it might perhaps be possible to use it in thousands of years, when observations are ten or a hundred times as accurate as they are now. To this I now reply

\SourceLeaf{490}
III) \textsc{Rohde}’s method is wrong even in theory, leads to unavoidable circular reasoning, and is consequently completely useless.

I have already pointed out that to find a planet’s distance from observations, one needs the distance of the Sun from the Earth at the same time, and that the more accurately the former is to be found, the more accurately the latter must be known. This requires knowing to the greatest accuracy both the elements of the Sun’s motion and all the perturbations that affect the distance of the Sun from the Earth, in particular those caused by the very planet whose distance is to be determined. Thus (and this is the \emph{point} of the matter) its mass must already be known. Otherwise the perturbation equations for the distance of the Sun from the Earth are incomplete, that distance is wrong, the calculated distance of the planet from the Sun is wrong, and the mean distance based on one or more true distances is wrong. An example will make this clear.

Suppose someone wanted to determine the mass of Mars by \textsc{Rohde}’s method. To determine its mean distance from observations, that person would need one or more true distances, which require the distances of the Sun from the Earth. In our older solar tables, however, all the equations that depend on Mars are missing because its mass is unknown (in Baron von \textsc{Zach}’s latest solar tables, Arg. IV. and Arg. XV). This omission alone could therefore introduce an error of $0.0000058 + 0.0000011 = 0.0000069$ in the Earth’s distance from the Sun. The resulting error in the distance of Mars from the Sun would be greater in the same proportion as this distance exceeds the former, and would thus amount to more than $0.00001$. It must be left to Captain \textsc{Rohde} himself to calculate what effect such an error in the distance would have on the mass. I also say nothing of the fact that to find the mass of a planet by \textsc{Rohde}’s method one would already need its own mass and those of all the other planets. Otherwise one could have neither an accurate theory of the Earth (to calculate the distance of the Sun from the Earth) nor an accurate theory of the planet itself. Its heliocentric longitude would have to be taken from the tables, and its actual distances could be used to find the mean distance only if they had first been freed from the effects of all external perturbations and thus made purely elliptical. But enough, and already far too much, about this immature work.
